%%%%%%%%%%%%%%%%%%%%%%%%%%%%%%%%%%%%%%%%%
% Beamer Presentation
% LaTeX Template
% Version 1.0 (10/11/12)
%
% This template has been downloaded from:
% http://www.LaTeXTemplates.com
%
% License:
% CC BY-NC-SA 3.0 (http://creativecommons.org/licenses/by-nc-sa/3.0/)
%
%%%%%%%%%%%%%%%%%%%%%%%%%%%%%%%%%%%%%%%%%

%----------------------------------------------------------------------------------------
%	PACKAGES AND THEMES
%----------------------------------------------------------------------------------------

\documentclass[xcolor=dvipsnames,french,graphics,aspectratio=169]{beamer}
\graphicspath{{images},{./graphics/fully_conv_nets}}

\mode<presentation> {

% The Beamer class comes with a number of default slide themes
% which change the colors and layouts of slides. Below this is a list
% of all the themes, uncomment each in turn to see what they look like.
\usepackage{amsmath,amssymb}

\usepackage{array}
\newcolumntype{L}[1]{>{\raggedright\let\newline\\\arraybackslash\hspace{0pt}}m{#1}}
\newcolumntype{C}[1]{>{\centering\let\newline\\\arraybackslash\hspace{0pt}}m{#1}}
\newcolumntype{R}[1]{>{\raggedleft\let\newline\\\arraybackslash\hspace{0pt}}m{#1}}

\usepackage{xcolor}
\usepackage{fancyvrb}
\usepackage{alltt}

\usepackage[absolute,showboxes,overlay]{textpos}
%\usepackage[texcoord,grid,gridunit=mm,gridcolor=red!10,subgridcolor=green!10]{eso-pic}%GRI!!!!

\textblockorigin{0pt}{0pt}                           %origine des positions (coin gauche) attention, axe y vers le bas !!
%\TPshowboxestrue                                  % affiche le contour des textblock
\TPshowboxesfalse                                % n'affiche pas le contour des textblock

\setlength{\TPHorizModule}{\paperwidth}\setlength{\TPVertModule}{\paperheight}
%\usetheme{default}
%\usetheme{AnnArbor}
%\usetheme{Antibes}
%\usetheme{Bergen}
%\usetheme{Berkeley}
%\usetheme{Berlin}
%\usetheme{Boadilla}
%\usetheme{CambridgeUS}
%\usetheme{Copenhagen}
%\usetheme{Darmstadt}
%\usetheme{Dresden}
%\usetheme{Frankfurt}
%\usetheme{Goettingen}
%\usetheme{Hannover}
%\usetheme{Ilmenau}
%\usetheme{JuanLesPins}
%\usetheme{Luebeck}
%\usetheme{Madrid}
%\usetheme{Malmoe}
%\usetheme{Marburg}
%\usetheme{Montpellier}
%\usetheme{PaloAlto}
%\usetheme{Pittsburgh}
%\usetheme{Rochester}
%\usetheme{Singapore}
%\usetheme{Szeged}
\usetheme{Warsaw}
\setbeamertemplate{headline}{%
\leavevmode%
  \hbox{%
    \begin{beamercolorbox}[wd=\paperwidth,ht=2.5ex,dp=1.125ex]{palette secondary}%quaternary
    \insertsectionnavigationhorizontal{\paperwidth}{}{\hskip0pt plus1filll}
    \end{beamercolorbox}%
  }
}

\defbeamertemplate*{footline}{shadow theme}
{%
  \leavevmode%
  \hbox{\begin{beamercolorbox}[wd=.5\paperwidth,ht=2.5ex,dp=1.125ex,leftskip=.3cm plus1fil,rightskip=.3cm]{author in head/foot}%
    \usebeamerfont{author in head/foot}\hfill\insertshortauthor
  \end{beamercolorbox}%
  \begin{beamercolorbox}[wd=.5\paperwidth,ht=2.5ex,dp=1.125ex,leftskip=.3cm,rightskip=.3cm plus1fil]{title in head/foot}%
    \usebeamerfont{title in head/foot}\insertshorttitle\hfill\insertframenumber\,/\,\inserttotalframenumber%
  \end{beamercolorbox}}%
  \vskip0pt%
}

\definecolor{mygray}{rgb}{0.3, 0.3, 0.3}
% --------------------------------------------------
%\usepackage[footheight=1em]{beamerthemeboxes}
%\addfootboxtemplate{\color{black}}{\color{white}
%     \hfill\insertframenumber/\inserttotalframenumber\hspace{2em}\null}
%\addtobeamertemplate{footline}{\insertframenumber/\inserttotalframenumber}
% --------------------------------------------------
% As well as themes, the Beamer class has a number of color themes
% for any slide theme. Uncomment each of these in turn to see how it
% changes the colors of your current slide theme.

%\usecolortheme{albatross}
%\usecolortheme{beaver}
%\usecolortheme{beetle}
%\usecolortheme{crane}%jaune
%\usecolortheme{dolphin}
%\usecolortheme{dove}% blanc sur blanc
%\usecolortheme{fly}%gris
%\usecolortheme{lily}% se parece a beamer
%\usecolortheme{orchid}
%\usecolortheme{rose}
%\usecolortheme{seagull}
%\usecolortheme{seahorse}
%\usecolortheme{whale}
%\usecolortheme{wolverine}

%\setbeamertemplate{footline} % To remove the footer line in all slides uncomment this line
%\setbeamertemplate{footline}[page number] % To replace the footer line in all slides with a simple slide count uncomment this line

\setbeamertemplate{navigation symbols}{} % To remove the navigation symbols from the bottom of all slides uncomment this line
}% mode presentation

\usepackage{graphicx} % Allows including images
\usepackage[export]{adjustbox}

\usepackage{booktabs} % Allows the use of \toprule, \midrule and \bottomrule in tables
\usepackage[latin1]{inputenc}
\usepackage{babel}
\usepackage{multimedia}
\usepackage{media9}
\usepackage{animate}% animategraphics
\usepackage{hyperref}
\usepackage{colortbl}  %Use to unconver columns in a table
\usepackage{tikz}

%\usepackage{enumitem}
\usepackage{array,arydshln}% http://ctan.org/pkg/array and \usepackage{} % dashlines

\newcolumntype{M}{>{\centering\arraybackslash}m{\dimexpr.25\linewidth-2\tabcolsep}}

\newenvironment<>{varblock}[2][.9\textwidth]{%
  \setlength{\textwidth}{#1}
  \begin{actionenv}#3%
    \def\insertblocktitle{#2}%
    \par%
    \usebeamertemplate{block begin}}
  {\par%
    \usebeamertemplate{block end}%
  \end{actionenv}}

%----------------------------------------------------------------------------------------
%	TITLE PAGE
%----------------------------------------------------------------------------------------

\title[ September 2025]{Semantic segmentation with U-net} % The short title appears at the bottom of every slide, the full title is only on the title page

\author{Beatriz Marcotegui} % Your name
\institute[MINES ParisTech] % Your institution as it will appear on the bottom of every slide, may be shorthand to save space
{
  Centre de Morphologie Math�matique \\ % Your institution for the title page
MINES ParisTech\\
PSL Research University\\
\medskip
\textit{beatriz.marcotegui@minesparis.psl.eu}\\ % Your email address
\bigskip

}
\date{September 2025} % Date, can be changed to a custom date



\AtBeginSection[] % Do nothing for \section*
{
\begin{frame}<beamer>
\frametitle{Plan}
\tableofcontents[currentsection]
\end{frame}
}




% UNCOMMENT NEXT LINES IN ORDER TO ADD TITLE FRAMES BETWEEN SECTIONS
% Add support for \subsubsectionpage
\def\subsubsectionname{\translate{Subsubsection}}
\def\insertsubsubsectionnumber{\arabic{subsubsection}}
\setbeamertemplate{subsubsection page}
{
  \begin{centering}
    {\usebeamerfont{subsubsection name}\usebeamercolor[fg]{subsubsection name}\subsubsectionname~\insertsubsubsectionnumber}
    \vskip1em\par
    \begin{beamercolorbox}[sep=4pt,center]{part title}
      \usebeamerfont{subsubsection title}\insertsubsubsection\par
    \end{beamercolorbox}
  \end{centering}
}
\def\subsubsectionpage{\usebeamertemplate*{subsubsection page}}

\AtBeginSection{\frame{\sectionpage}}
\AtBeginSubsection{\frame{\subsectionpage}}

\newcommand{\dl}{\delta}
\newcommand{\source}[1]{\begin{textblock*}{4cm}(8.7cm,8.6cm)
    \begin{beamercolorbox}[ht=0.5cm,right]{framesource}
      \usebeamerfont{framesource}\usebeamercolor[fg]{framesource} Credits: {#1}
    \end{beamercolorbox}
\end{textblock*}}

\newcommand{\x}{\mathbf{x}}
\newcommand{\y}{\mathbf{y}}
\newcommand{\X}{\mathbf{X}}
\newcommand{\Y}{\mathbf{Y}}
\newcommand{\W}{\mathbf{W}} %Weight
\newcommand{\bias}{\mathbf{b}}%Bias
\newcommand{\act}{\texttt{g}}%Activation
\newcommand{\lr}{\eta} % learning rate
\newcommand{\loss}{L}
\newcommand{\pdata}{\hat{p}_{\texttt{data}}}
\newcommand{\nsize}{N}
\newcommand{\nfeatures}{P}
\newcommand{\param}{\boldsymbol{\theta}}
\newcommand{\featmap}{\boldsymbol{\phi}}
\newcommand{\EV}{\mathbb{E}}

\newcommand{\R}{\mathbb{R}}
\newcommand{\N}{\mathbb{N}}
\newcommand{\Z}{\mathbb{Z}}
\newcommand{\myvec}[1]{\ensuremath{\begin{pmatrix}#1\end{pmatrix}}}

%%%%%%%%%%%%%%%%%%%%%%%%%%%%%%%%%%%%%%%%%%%%%%%%%% ``ADD VIDEO'' DEFINITION
\newcommand{\addvideo}[3]{
% Macro to add video :
% Argument 1 : video source location
% Argument 2 : image preview file location
% Argument 3 : width of video
\begin{center}
    \includemedia
    [
        width  = #3\linewidth,
        height  = #3\linewidth,
        activate=pageopen,
        addresource = #1,               %Inclut le fichier source dans le fichier pdf
        flashvars = {source=#1 & scaleMode = letterbox&autoPlay=true&loop=true}     % Appelle le fichier source dans le lecteur
    ]
    {#2}{VPlayer.swf} % Arg #2 is used as preview image
\end{center}
}

\begin{document}
\tikzstyle{RectObject}=[rectangle,fill=white,draw,line width=0.5mm]
\tikzstyle{RectObject2}=[rectangle,fill=white,draw,line width=0.7mm color = <blue>]
\tikzstyle{line}=[draw]
\tikzstyle{arrow}=[draw, -latex]
\setlength{\unitlength}{1cm}

%%%%%%%%%%%%%%%%%%%%%%%%%%%%%%%%%%%%%%%%%%%%%%%%%%%%%%%%%%%%
\setbeamercovered{transparent=30}

\begin{frame}
\titlepage % Print the title page as the first slide
\end{frame}

\begin{frame}
\frametitle{Plan } % Table of contents slide, comment this block out to remove it
\tableofcontents[currentsection] % Throughout your presentation, if you choose to use \section{} and \subsection{} commands, these will automatically be printed on this slide as an overview of your presentation
\end{frame}


\section{Introduction}
%%%%%%%%%%%%%%%%%%%%%%%%%%%%%%%%%%%%%%%%%%%%%%%%%%%%%%%%%%%%
    \setbeamercovered{invisible}


    %%%%%%%%%%%%%%%%%%%%%%%%%%%%%%%%%%%%%%%%%%%%%%%%%%%%%%%%%%%%
\begin{frame}{Some vocabulary}

  %..............................
  \begin{itemize}{}

  \item \textbf{Semantic segmentation}: each pixel gets a "class" label
\onslide<4->{  \item \textbf{Instance segmentation}:  each pixel takes an "instance" label.}
\onslide<5->{  \item \textbf{Panoptic segmentation}: each pixel takes "class" and "instance" labels.}
  \end{itemize}
  %..............................

%  Note: \textbf{panoptic} definition (Larousse): Said of a building (penitentiary, hospital, etc.) from which, from an internal observation deck, one can see the entire interior.
\begin{small}
\begin{center}
  %::::::::::::::::::::::::::::::
  \begin{tabular}{cccc}

      \onslide<2->{original}&\onslide<3->{\textbf{semantic seg}}&\onslide<4->{\textbf{instance seg}}&\onslide<5->{\textbf{panoptic seg}}\\

      \onslide<2->{\includegraphics[width=0.22\textwidth]{../images/berkeley_elephants.png}}&
      \onslide<3->{
\includegraphics[width=0.22\textwidth]{../images/berkeley_semantic_seg.png}}&
      \onslide<4->{\includegraphics[width=0.22\textwidth]{../images/berkeley_elephants_instance.png}}&
      \onslide<5->{
\includegraphics[width=0.22\textwidth]{../images/berkeley_semantic_seg.png}}\\
      & & &      \onslide<5->{\includegraphics[width=0.22\textwidth]{../images/berkeley_elephants_instance.png}}

  \end{tabular}
  %::::::::::::::::::::::::::::::
\end{center}
\end{small}

%% %::::::::::::::::::::::::::::::
%% \begin{tabular}{ccc}
%% \includegraphics[width=0.3\textwidth]{../images/hierarchy/car_ori.png}&
%% \includegraphics[width=0.3\textwidth]{../images/hierarchy/car_ws.png}&
%% \includegraphics[width=0.3\textwidth]{../images/hierarchy/car_hierarchy.png}\\

%% \end{tabular}
%% %::::::::::::::::::::::::::::::

\end{frame}


%% %%%%%%%%%%%%%%%%%%%%%%%%%%%%%%%%%%%%%%%%%%%%%%%%%%%%%%%%%%%%
%% \begin{frame}{Pascal visual object classes segmentation challenge 2012 \cite{everingham_pascal_2014}}


%%   \begin{itemize}
%%   \item 1464 training and 1449 validation images
%%   \item automatic online test, with unknown images
%%   \item 20 image categories (cat, sofa, motorbike, person, etc.)
%%   \end{itemize}

%%   \begin{figure}
%%     \includegraphics[height=2.5cm]{../images/pascal_cat.jpg}
%%     \includegraphics[height=2.5cm]{../images/pascal_cat_seg.png}\\
%%     \includegraphics[height=2.5cm]{../images/pascal_moto.jpg}
%%     \includegraphics[height=2.5cm]{../images/pascal_moto_seg.png}
%%   \end{figure}

%% \end{frame}
\section{Metrics}

%%%%%%%%%%%%%%%%%%%%%%%%%%%%%%%%%%%%%%%%%%%%%%%%%%

\begin{frame}{Measures used in image processing}


  %::::::::::::::::::::::::::::::
  \begin{columns}
    \begin{column}{0.6\textwidth}
      Let $A$ and $B$ be two sets, not simultaneously empty ({\begin{small} Credits: public domain wikipedia.org. \end{small}})

    \end{column}
    \begin{column}{0.3\textwidth}
      \includegraphics[width=0.8\textwidth]{../images/Intersection_of_sets_A_and_B.png}
    \end{column}
  \end{columns}
  %::::::::::::::::::::::::::::::


  %::::::::::::::::::::::::::::::
  \begin{columns}
    \begin{column}{0.5\textwidth}
      \begin{block}{Jaccard index}
        \[
        J(A,B) = \frac{|A \cap B|}{|A \cup B|}
        \]
      \end{block}

    \end{column}
    \begin{column}{0.5\textwidth}
      \begin{block}{Dice coefficient}
        \[
        D(A,B) = \frac{2|A \cap B|}{|A|+|B|}
        \]
      \end{block}

    \end{column}
  \end{columns}
  %::::::::::::::::::::::::::::::

  \begin{block}{Properties}
    \begin{itemize}
    \item $\forall A, B: 0 \leq J(A, B) \leq D(A,B) \leq 1$
    \item If $A=B$, then $D(A,B) = J(A,B) = 1$
    \item If $ A \cap B = \emptyset$, then $D(A,B) = J(A,B) = 0$
    \end{itemize}
  \end{block}


\end{frame}


%%%%%%%%%%%%%%%%%%%%%%%%%%%%%%%%%%%%%%%%%%%%%%%%%%

\begin{frame}{Generalization to $[0,1]$}

  \begin{block}{Jaccard index}
    Also named Intersection over Union (IoU)
        \[
        J(A,B) = \frac{|A \cap B|}{|A \cup B|}
        \]
      \end{block}

      But $\y$ and $\hat{\y}$ are in $[0,1]^n \ldots$

\pause

  %% \begin{block}{Dice similarity}
  %%   \[
  %%   D(\y, \hat{\y}) = \frac{2\sum_i y_i\hat{y}_i}{\sum_i y_i + \sum_i \hat{y}_i}
  %%   \]
  %% \end{block}

  \begin{block}{Jaccard similarity}
    \[
    J(\y, \hat{\y}) = \frac{\sum_i y_i\hat{y}_i}{\sum_i y_i + \sum_i \hat{y}_i - \sum_i y_i\hat{y}_i}
    \]
  \end{block}

($\y$ and $\hat{\y}$ are not simultaneously equal to $0$)

\end{frame}

%%%%%%%%%%%%%%%%%%%%%%%%%%%%%%%%%%%%%%%%%%%%%%%%%%

\begin{frame}{Corresponding loss function}


  %% \begin{block}{Dice loss}
  %%   \[
  %%   d(\y, \hat{\y}) = 1-\frac{2\sum_i y_i\hat{y}_i}{\sum_i y_i + \sum_i \hat{y}_i}
  %%   \]
  %% \end{block}

  \begin{block}{Jaccard loss}
    \[
    j(\y, \hat{\y}) = \textcolor{blue}{1-}\frac{\sum_i y_i\hat{y}_i}{\sum_i y_i + \sum_i \hat{y}_i - \sum_i y_i\hat{y}_i + \epsilon}
    \]
  \end{block}

    Constant $\epsilon$, which is typically ``small'', keeps the denominator ``far enough'' from zero.


\end{frame}

%%%%%%%%%%%%%%%%%%%%%%%%%%%%%%%%%%%%%%%%%%%%%%%%%%
%%%%%%%%%%%%%%%%%%%%%%%%%%%%%%%%%%%%%%%%%%%%%%%%%%%%%%%%%%%%
\begin{frame}[fragile]{Jaccard Loss code (keras)}
\begin{verbatim}
def jaccard1_coef(y_true, y_pred, smooth=SMOOTH):
    y_true_f = K.flatten(y_true)
    y_pred_f = K.flatten(y_pred)
    intersection = K.sum(y_true_f * y_pred_f)
    union = K.sum(y_true_f) + K.sum(y_pred_f) - intersection
    return (intersection + smooth) / (union + smooth)


def jaccard1_loss(y_true, y_pred, smooth=SMOOTH):
    return 1 - jaccard1_coef(y_true, y_pred, smooth=SMOOTH)
\end{verbatim}
  \end{frame}
\begin{frame}{Jaccard loss function - variant}


  %% \begin{block}{Dice loss}
  %%   \[
  %%   d(\y, \hat{\y}) = 1-\frac{2\sum_i y_i\hat{y}_i}{\sum_i y_i^2 + \sum_i \hat{y}_i^2 + \epsilon}
  %%   \]
  %% \end{block}

  \begin{block}{Jaccard2 loss}
    \[
    j_2(\y, \hat{\y}) = 1-\frac{\sum_i y_i\hat{y}_i}{\sum_i y_i^2 + \sum_i \hat{y}_i^2 - \sum_i y_i\hat{y}_i + \epsilon}
    \]
  \end{block}

  This version works slightly better than standard Jaccard loss ($j_1$).

  \vspace*{0.5cm}
  \parbox[c]{0.9\textwidth}{
\begin{scriptsize}
    %..............................
    \begin{itemize}{}
    \item Cha, S. H. (2007). Comprehensive survey on distance/similarity measures between probability density functions. City, 1(2), 1.

    \item
 Duque-Arias, D. et al. (2021). On power Jaccard losses for semantic segmentation. In VISAPP 2021.
      \end{itemize}
    %..............................
\end{scriptsize}
}
\end{frame}




%%%%%%%%%%%%%%%%%%%%%%%%%%%%%%%%%%%%%%%%%%%%%%%%%%%%%%%%%%%%
\begin{frame}{Loss functions}
    \begin{itemize}
    \item Cross-entropy: $ -\sum\limits_i y_i\log(\hat{y}_i)$
    \end{itemize}

\end{frame}

%%%%%%%%%%%%%%%%%%%%%%%%%%%%%%%%%%%%%%%%%%%%%%%%%%%%%%%%%%%%

\begin{frame}{Cross entropy}
  Classification loss function.


  It is the sum of the true probability ($y_i$) multiplied by the log of the prediction ($\hat{y}_i$) over all classes of the distribution.

\textcolor{blue}{  \centerline{  Cross-entropy:  $CE=-\sum\limits_i y_i\log(\hat{y}_i)$}}

  %::::::::::::::::::::::::::::::
  \begin{tabular}{ccccc}
    input image&
    {\large $\boldsymbol{\Rightarrow}$}&
    MODEL &
    {\large $\boldsymbol{\Rightarrow}$}&
    vector of probabilities
  \end{tabular}
  %::::::::::::::::::::::::::::::
  \centerline{  \includegraphics[width=0.2\textwidth]{../images/dog_image.png}}

  %::::::::::::::::::::::::::::::
  \begin{center}
  %::::::::::::::::::::::::::::::
  \begin{tabular}{|c|c|c|c|c|}

    \hline
    & dog&cat&elephant&snike\\
    \hline
    True probability & 1 & 0 & 0 & 0\\
    \hline
    Predicted prob.  &0.91&0.08&0.01&0\\
    \hline
  \end{tabular}
  \end{center}


Usually softmax is found at the end of DL classification systems, in order to convert output scores in probabilities (between [0,1] and sum to 1).

\end{frame}

%%%%%%%%%%%%%%%%%%%%%%%%%%%%%%%%%%%%%%%%%%%%%%%%%%%%%%%%%%%%
\begin{frame}{Softmax operator}

  \begin{block}{Definition}
    The softmax operator $\sigma: \R^d \rightarrow \R^d$ is given by:
    \[
    \forall \mathbf{x}  \in \R^d, \forall k \in \{1, \ldots, d\}:\quad \sigma(\mathbf{x})_k = \frac{e^{\mathbf{x}_k}}{\sum\limits_{i=1}^d e^{\mathbf{x}_i}}
    \]
  \end{block}

  \pause

  \begin{columns}
    \begin{column}{.3\textwidth}
      \begin{block}{Some properties}
        \begin{itemize}
        \item $0 < \sigma(\mathbf{x})_k < 1$
        \item $\sum\limits_{i=1}^d \sigma(\mathbf{x})_i = 1$
        \end{itemize}
      \end{block}

    \end{column}

    \pause

    \begin{column}{.7\textwidth}

      \begin{block}{Example}
        \[
        \mathbf{x} = \myvec{10.1\\0\\-4.3\\1.33}  \quad \sigma(\mathbf{x}) \approx \myvec{0.9998\\0.000041\\0.00000056\\0.00016}
        \]

      \end{block}



    \end{column}
  \end{columns}



\end{frame}


%%%%%%%%%%%%%%%%%%%%%%%%%%%%%%%%%%%%%%%%%%%%%%%%%%%%%%%%%%%%
\begin{frame}{Cross entropy}
  \begin{block}{  Kullback - Leibler divergence}
    The dissimilarity between two probability distributions is estimated by the KL divergence.
  \end{block}
  %..........

  Consider two probability distributions $P$ and $Q$: the true probability ($P=y$) and  the predicted probability ($Q=\hat{y}_i$).

  $D_{KL}(P,Q)=\sum\limits_i P_i\log \frac{P_i}{Q_i} = \sum\limits_i P_i [\log (P_i)-\log(Q_i)]= \sum\limits_i P_i \log(P_i) - \sum\limits_i P_i \log(Q_i) $

  The first term $\sum\limits_i P_i \log(P_i)$ is the entropy of the ground truth labels, that does not depend on $\theta$. Therefore:

  $ \underset{\theta}{\mathrm{argmin}}[ D_{KL}(P,Q)] = \underset{\theta}{\mathrm{argmin}}[- \sum\limits_i P_i \log(Q_i)] \Rightarrow \textcolor{blue}{\text{Cross Entropy formula}}$
%https://www.youtube.com/watch?v=SxGYPqCgJWM
\end{frame}

%%%%%%%%%%%%%%%%%%%%%%%%%%%%%%%%%%%%%%%%%%%%%%%%%%%%%%%%%%%%
\begin{frame}{Cross entropy}
  %::::::::::::::::::::::::::::::
  \begin{center}

 \includegraphics<1>[width=0.6\textwidth]{../images/cross_entropy_loss_1.png}
 \includegraphics<2>[width=0.6\textwidth]{../images/sq_CE_losses_1.png}
\includegraphics<3>[width=0.6\textwidth]{../images/sq_CE_grad_1.png}
  \end{center}
  %::::::::::::::::::::::::::::::

\only<3>{ Cross entropy: higher gradient for worst predictions. $\Rightarrow$ faster convergence}
\end{frame}

%%%%%%%%%%%%%%%%%%%%%%%%%%%%%%%%%%%%%%%%%%%%%%%%%%%%%%%%%%%%
\begin{frame}{Some losses comparison}
  \centerline{ \includegraphics[width=0.6\textwidth]{../images/jac_sq_CE_losses_1_limit0-1.png}}

  Sometimes, several losses are combined.
  \end{frame}
%%%%%%%%%%%%%%%%%%%%%%%%%%%%%%%%%%%%%%%%%%%%%%%%%%%%%%%%%%%%

\begin{frame}{Focal loss}
  %..........
  \begin{block}{Goal}
Designed to fight against class imbalance, focuses training on a
sparse set of "hard" examples and prevents the vast number of "easy"
samples from overwhelming the detector.
    \end{block}
      $$
      CE(p,y)=
    \left\lbrace
      \begin{array}{cc}
        -log(p)  & \text{if } y = 1\\
        -log(1-p) & otherwise
      \end{array}\right.
      $$

      y: ground truth class

      $p \in [0,1]$: estimated probability for the class y = 1

            $$
      p_t=
    \left\lbrace
      \begin{array}{cc}
        p  & \text{if } y = 1\\
        (1-p) & otherwise
      \end{array}\right.
      $$

      $$ CE(p,y)=CE(p_t)=-log(p_t)$$

      {\scriptsize   Focal Loss for Dense Object Detection
  https://arxiv.org/pdf/1708.02002v2.pdf
}
\end{frame}
%%%%%%%%%%%%%%%%%%%%%%%%%%%%%%%%%%%%%%%%%%%%%%%%%%%%%%%%%%%%
\begin{frame}{Focal loss}
  %..........
  \begin{block}{definition}
  $$FL(p_t)=-(1-p_t)^\gamma log(p_t)$$
    \end{block}
\centerline{  \includegraphics[width=0.6\textwidth]{../images/focal_loss.png}}

\end{frame}
%%%%%%%%%%%%%%%%%%%%%%%%%%%%%%%%%%%%%%%%%%%%%%%%%%%%%%%%%%%%
\begin{frame}{Multi-class segmentation}
  \parbox{0.9\textwidth}{Each class is set in a different output channel (one-hot encoding). Then binary segmentation is applied.}

  \centerline{  \includegraphics[width=0.7\textwidth]{../images/categorical_2D.png}}


  \end{frame}

%%%%%%%%%%%%%%%%%%%%%%%%%%%%%%%%%%%%%%%%%%%%%%%%%%%%%%%%%%%%
\begin{frame}{Multi-class segmentation}
  Each class is set in a different output channel.

  For example:

  %::::::::::::::::::::::::::::::
  \begin{columns}

    \begin{column}{0.15\textwidth}
      \begin{center}
      %::::::::::::::::::::::::::::::
      \begin{tabular}{|c|}
        \hline
        \textcolor{green}{1}\\
        \hline
        \textcolor{red}{0}\\
        \hline
        \textcolor{blue}{2}\\
        \hline
        \textcolor{blue}{2}\\
        \hline
        \textcolor{green}{1}\\
        \hline
        \textcolor{red}{0}\\
        \hline
        \textcolor{blue}{2}\\
        \hline
        \textcolor{blue}{2}\\
        \hline
      \end{tabular}
      %::::::::::::::::::::::::::::::
          \end{center}
    \end{column}

    \begin{column}{0.1\textwidth}
      {\large $$\boldsymbol{\Rightarrow}$$}
    \end{column}

    \begin{column}{0.15\textwidth}
      \begin{center}
      \begin{tabular}{|c|c|c|}
        \hline
        0&\textcolor{green}{1}&0\\
        \hline
        \textcolor{red}{1}&0 &0\\
        \hline
        0&0&\textcolor{blue}{1}\\
        \hline
        0 &0&\textcolor{blue}{1}\\
        \hline
        0 &\textcolor{green}{1} &0\\
        \hline
        \textcolor{red}{1}&0 &0\\
        \hline
        0 &0&\textcolor{blue}{1}\\
        \hline
        0 &0&\textcolor{blue}{1}\\
        \hline
      \end{tabular}
      %::::::::::::::::::::::::::::::
    \end{center}
    \end{column}

  \end{columns}

  $$Y2= tf.keras.utils.to\_categorical(Y)$$
  %::::::::::::::::::::::::::::::
\end{frame}
%%%%%%%%%%%%%%%%%%%%%%%%%%%%%%%%%%%%

%%%%%%%%%%%%%%%%%%%%%%%%%%%%%%%%%%%%%%%%%%%%%%%%%%%%%%%%%%%%
\begin{frame}{multi-class segmentation metric}
  %..........
  \begin{block}{mean Intersection over Union (mIoU)}
  Usual multi-class segmentation metric:

$$  mIoU = \frac{1}{N} \sum_{i=1} ^N IoU_i$$

  $N$: number of classes, $IoU_i$: IoU of class i.
    \end{block}
  \end{frame}
\section{Optimizers}

%%%%%%%%%%%%%%%%%%%%%%%%%%%%%%%%%%%%%%%%%%%%%%%%%%%%%%%%%%%%
\begin{frame}{Optimizers}
%https://www.youtube.com/watch?v=YnQJTfbwBM8&list=PL5-TkQAfAZFbzxjBHtzdVCWE0Zbhomg7r&index=6
  Method used to update the NN weights so that the loss in the training set is reduced.

  The problem is strongly non-convex. We do not have convergence guarantees. Only a few guidelines.
%..............................
  \begin{itemize}{}
  \item Gradient descent
  \item Stochastic Gradient Descent (SGD)


\end{itemize}
%..............................

  \end{frame}
%%%%%%%%%%%%%%%%%%%%%%%%%%%%%%%%%%%%%%%%%%%%%%%%%%%%%%%%%%%%
\begin{frame}[fragile]{Gradient descent}
  \begin{alltt}
    for t in nb_steps:
        dw = compute_gradient(lossF,training_samples,w)
        w = w - lr * dw
    \end{alltt}
  %::::::::::::::::::::::::::::::
  \begin{columns}
    \begin{column}{0.5\textwidth}
      \centerline{  \includegraphics[width=0.8\textwidth]{../images/gradient_descent.jpeg}}
      {\scriptsize \parbox {0.95\textwidth}{ From https://towardsdatascience.com/linear-regression-using-gradient-descent-97a6c8700931}}
    \end{column}
    \begin{column}{0.5\textwidth}
      $ dw = \bigtriangledown_\theta \sum_i^N  L(f(x_i,\theta_{t-1}),y_i)$

      with $(x_i,y_i)$, the training set and N the number of training samples.

    \end{column}
  \end{columns}
  %::::::::::::::::::::::::::::::

  \end{frame}
%%%%%%%%%%%%%%%%%%%%%%%%%%%%%%%%%%%%%%%%%%%%%%%%%%%%%%%%%%%%
\begin{frame}[fragile]{Stochastic Gradient Descent}
  \begin{alltt}
    for t in nb_steps:
       \textcolor{blue}{SelectBatch}
       dw = compute_gradient(lossF,\textcolor{blue}{batch},w)
       w = w - lr * dw
  \end{alltt}

  $ dw = \bigtriangledown_\theta \sum_i^m  L(f(x_i,\theta_{t-1}),y_i)$

  with $(x_i,y_i)$, the training set and m the number of batch samples.

  %..........
  \begin{block}{Hyper-parameters}
    %..............................
    \begin{itemize}{}
    \item weight initialization
    \item learning rate
    \item number of iterations
    \item batch size (stochastic!) because it is computationaly expensive to compute the loss on the whole dataset
    \end{itemize}
    %..............................

    \end{block}
  \end{frame}
%%%%%%%%%%%%%%%%%%%%%%%%%%%%%%%%%%%%%%%%%%%%%%%%%%%%%%%%%%%%


\begin{frame}[fragile]{SGD + momentum}

  %..........
  \begin{block}{Momentum}
     Accelerates optimization, integrating gradient with an exponentially decaying moving average


    \end{block}

  \begin{alltt}
    v = 0
    for i in nb_steps:
        dw = compute_gradient(lossF,data,w)
        \textcolor{blue}{v = rho * v + dw}
        w = w - lr * \textcolor{blue}{v}
  \end{alltt}
  %..............................
  \begin{itemize}{}
  \item Intuition: a ball going down a hill. The ball does not follow the gradient direction because of its velocity.
  \item $\rho$ typically 0.9 or 0.99
  \end{itemize}
  %..............................


  %% \begin{alltt} EQUIVALENT
  %%   v = 0
  %%   for i in nb_steps:
  %%       dw = compute_gradient(lossF,data,w)
  %%       v = rho * v  - lr * dw
  %%       $ w^{n+1} = w^{n} + v$

  %%   \end{alltt}

\end{frame}

%%%%%%%%%%%%%%%%%%%%%%%%%%%%%%%%%%%%%%%%%%%%%%%%%%%%%%%%%%%%


\begin{frame}[fragile]{SGD + momentum}

  $$ v_i = \rho v_{i-1} + dw_i$$
    $ v_0 = dw_0$

  $v_1 = \rho * dw_0 + dw_1$

  $v_2 = \rho * v_1 + dw_2 = \rho (\rho * v_0 + dw_1)+ dw_2 = \rho^2 dw_0 + \rho dw_1 + dw_2$

  $...$

  $ v_i =\rho^i dw_0 + \rho^{i-1} dw_1 + ....+  dw_i $

\textcolor{blue}{    $$ v_i =\rho^i dw_0 + \rho^{i-1} dw_1 + ....+  dw_i $$}

%..............................
\begin{itemize}{}

    \item If gradient is in the same direction it is amplified

    \item Abrupt changes are smoothed


\end{itemize}
%..............................

{\scriptsize Sutskever et al."On the importance of initialization and momentun in DL" ICML 2013.}
\end{frame}

%%%%%%%%%%%%%%%%%%%%%%%%%%%%%%%%%%%%%%%%%%%%%%%%%%%%%%%%%%%%

\begin{frame}{SGD difficulties}


  %..............................
  \begin{itemize}{}
  \item  If the gradient is high in one dimension and low in another one:

  %..............................
    \begin{itemize}{}
    \item The optimization process may oscillate ("zigzag") if the learning rate is large
    \item The optimization process may be slow, if the learning rate is small

      In both cases, too many iterarations are required for convergence, if it converges!
    \end{itemize}

      %::::::::::::::::::::::::::::::
    \begin{center}
    \begin{columns}
      \begin{column}{0.5\textwidth}
        \centerline{\includegraphics<1>[width=0.6\textwidth]{../images/DL_optimization_path_01.png}}
      \end{column}
      \begin{column}{0.5\textwidth}
        \centerline{\includegraphics<1>[width=0.6\textwidth]{../images/DL_optimization_path_02.png}}
      \end{column}
    \end{columns}
    \end{center}
    %::::::::::::::::::::::::::::::

\onslide<2->{    \item local minima}
\onslide<3->{    \item saddle points: increases in one dimension and decreases in another one}
\onslide<4->{    \item noisy "stochastic" estimations}

  \end{itemize}
  %..............................

  % having a high condition number. Can be estimated by the ratio between the singular value of the Hessian matrix.
  % saddle points have 0 gradient. Increases in one dimension and decreases in another one
\end{frame}

%%%%%%%%%%%%%%%%%%%%%%%%%%%%%%%%%%%%%%%%%%%%%%%%%%%%%%%%%%%%

%%%%%%%%%%%%%%%%%%%%%%%%%%%%%%%%%%%%%%%%%%%%%%%%%%%%%%%%%%%%
\begin{frame}[fragile]{AdaGrad}
\textcolor{blue}{  Adaptive} learning rates (per parameter).
  \begin{alltt}
    grad_sq = 0
    for i in nb_steps:
        dw = compute_gradient(lossF,data,w)
        \textcolor{blue}{grad_sq +=  dw * dw}
        w = w - lr*dw/(grad_sq+epsilon)
    \end{alltt}

  The squared gradient is always positive. Its sum is increasing over time. Then, AdaGrad decreases the learning rate over time.

  The square gradient may become too big. The optimization progress may stop.

  Duchi et al. "Adaptive subgradient methods for online learning and stochastic optimization" JMLR 2011.
\end{frame}

%%%%%%%%%%%%%%%%%%%%%%%%%%%%%%%%%%%%%%%%%%%%%%%%%%%%%%%%%%%%
\begin{frame}[fragile]{RMSProp: "Leak Adagrad"}
  \begin{alltt}
    grad_sq = 0
    for i in nb_steps:
        dw = compute_gradient(lossF,data,w)
        \textcolor{blue}{grad_sq = decay_rate * grad_sq + (1-decay_rate) * dw * dw}
        w = w - lr*dw/(grad_sq+epsilon)
    \end{alltt}

  Intuition: do not slow down over the optimization process.

  Momentum based optimization can overshoot and then come back.

  AdaGrad and RMSProp arrest progress along the fast moving direction while simultaneously accelerating progress along slow moving direction. $\Rightarrow$ Bends directly to the minimum.
\end{frame}

%%%%%%%%%%%%%%%%%%%%%%%%%%%%%%%%%%%%%%%%%%%%%%%%%%%%%%%%%%%%
\begin{frame}[fragile]{Adam: RMSProp + momentum}
  Bias at the beginning of the optimization process, gradient is close to 0.


  \begin{alltt}
    grad_sq = 0
    for i in nb_steps:
        dw = compute_gradient(lossF,data,w)
        moment1 = beta1 * moment1 + (1-beta1) * dw
        moment2 = beta2 * moment2 + (1-beta2) * dw * dw

        momment1_unbias = moment1/(1-beta1**i)
        momment2_unbias = moment2/(1-beta2**i)
        w = w - lr*moment1_unbias/(moment2_unbias+epsilon)

    \end{alltt}


default values: $ \beta_1=0.9; \beta_2 = 0.999$

  Kingma and Ba, "Adam: A method for stochastic optimization", ICLR 2015.
\end{frame}
%%%%%%%%%%%%%%%%%%%%%%%%%%%%%%%%%%%%%%%%%%%%%%%%%%%%%%%%%%%%
\begin{frame}{Weigth Decay}

{\bf  Idea}: keep weights small unless absolutely necessary. A regularization technique that prevents overfitting.

  %::::::::::::::::::::::::::::::
  \begin{columns}
    \begin{column}{0.35\textwidth}
        $$ \mathcal{L}_{total} = \mathcal{L}_{original}+ \lambda \sum w_i^2$$

    \end{column}
    \begin{column}{0.65\textwidth}
        \begin{itemize}{}
        \item

          $ \mathcal{L}_{original}$: Original loss, i.e. cross entropy, dice...
        \item $ \sum w_i^2$ : $L_2$ norm of all weights
        \item $\lambda$: weight decay coefficient
        \end{itemize}
    \end{column}


  \end{columns}
  %::::::::::::::::::::::::::::::
  %..............................

  %% \vspace*{0.5cm}

  %%   \centerline{      \includegraphics[width=0.95\textwidth]{weight_decay_graphics2.png}}

  %..............................

  \end{frame}

%%%%%%%%%%%%%%%%%%%%%%%%%%%%%%%%%%%%%%%%%%%%%%%%%%%%%%%%%%%%
\begin{frame}[fragile]{AdamW}

  Adam: weight decay added as L2 regularization term to loss function.

  AdamW: weight decay applied directly to weight updates

   \begin{alltt}
   w = w - lr (moment1_unbias /(moment2_unbias + epsilon) + lambda w)
   \end{alltt}
  %::::::::::::::::::::::::::::::
   \begin{tabular}{|c|c|c|}
     \hline
     & Adam&AdamW\\
    \hline
    Weight Decay   &  \parbox[m]{5cm}{L2 regularization applied to the loss}  & \parbox[t]{5cm}{Applied directly to \\weight updates }\\
    \hline

     Impact         & \parbox[c]{5cm}{Can interfere with adaptive lr} & \parbox[t]{5cm}{Decoupled, leads to \\consistent updates } \\
        \hline

        Generalization & \parbox[t]{5cm}{Can cause overfitting\\ in complex models}  & \parbox[t]{5cm}{Better regularization \\and generalization }\\
        \hline

  \end{tabular}
  %::::::::::::::::::::::::::::::


\end{frame}
%%%%%%%%%%%%%%%%%%%%%%%%%%%%%%%%%%%%%%%%%%%%%%%%%%%%%%%%%%%%
\begin{frame}{Adam vs AdamW}

  %..............................
  \begin{itemize}{}
  \item {\bf Adam}: may converge faster but can suffer from poor generalization
  \item {\bf AdamW}: in case of large models, AdamW may help to achieve better generalization and stable convergence.

  \end{itemize}
  %..............................


  \end{frame}
\section{From classification to image-to-image translation}


%%%%%%%%%%%%%%%%%%%%%%%%%%%%%%%%%%%%
%% \begin{frame}{VGG16: a network for image classification}

%%   \begin{center}
%%     \begin{figure}
%%       \centering
%%       \includegraphics[width=0.65\textwidth]{vgg16.png}
%%       \source{VGG16 (From https://www.cs.toronto.edu/~frossard/post/vgg16/)}
%%     \end{figure}
%%   \end{center}


%% \end{frame}



%%%%%%%%%%%%%%%%%%%%%%%%%%%%%%%%%%%%
\begin{frame}{From classification to image-to-image translation}
  %::::::::::::::::::::::::::::::
  \begin{columns}
    \begin{column}{0.5\textwidth}
      \begin{figure}
      \includegraphics[width=0.8\textwidth]{image_classif.png}
      \end{figure}

    \end{column}
          \pause
          \begin{column}{0.5\textwidth}
            \begin{figure}
      \includegraphics[width=0.8\textwidth]{image_transf.png}
      \end{figure}

    \end{column}

  \end{columns}
  %::::::::::::::::::::::::::::::





\end{frame}


%%%%%%%%%%%%%%%%%%%%%%%%%%%%%%%%%%%%
\begin{frame}{From classification nets to image-to-image nets}
  %::::::::::::::::::::::::::::::
  \begin{columns}
    \begin{column}{0.5\textwidth}
      \begin{figure}
      \includegraphics[width=0.5\textwidth]{image_transf1.png}
      \end{figure}

    \end{column}
          \pause
          \begin{column}{0.5\textwidth}
            \begin{figure}
      \includegraphics[width=0.5\textwidth]{image_transf2.png}
      \end{figure}

    \end{column}

  \end{columns}
  %::::::::::::::::::::::::::::::
\end{frame}

%%%%%%%%%%%%%%%%%%%%%%%%%%%%%%%%%%%%
\begin{frame}{From classification nets to image-to-image nets}
  %::::::::::::::::::::::::::::::
  \begin{columns}
    \begin{column}{0.5\textwidth}
      \begin{figure}
      \includegraphics[width=0.5\textwidth]{image_transf_fin.png}
      \end{figure}

    \end{column}
          \pause
          \begin{column}{0.5\textwidth}

      \begin{itemize}
      \item The neuron membrane segmentation challenge winner~\cite{ciresan_deep_2012} used this strategy. It is inefficient.
      \item The idea to compute the whole output image with a single pass through the network had already been proposed \cite{ning_toward_2005, jain_supervised_2007} (but with architectures more complex than those used today).
      \end{itemize}

    \end{column}

  \end{columns}
  %::::::::::::::::::::::::::::::



\end{frame}


%%%%%%%%%%%%%%%%%%%%%%%%%%%%%%%%%%%%%%%%%%%%%%%%%%
\frame{
  \frametitle{The simplest image-to-image architecture}

  \begin{block}{Example: plain CNN \cite{pang_cell_2010}}
    \begin{figure}
      \includegraphics[width=0.7\textwidth]{plain_convnet.png}
    \end{figure}
  \end{block}

  \pause

Note that today we would use ReLU activations on all layers except the last.

\begin{block}{}%quizzblock
  How many parameters does the first layer contain? The last? The second?
\end{block}%quizzblock

\pause

\mode<beamer>{$(9+1) \times d$ ;  $9d+1$ ; $(9d+1) \times d$}

}

%%%%%%%%%%%%%%%%%%%%%%%%%%%%%%%%%%%%
\begin{frame}<beamer>{Going back to the original image size}

    \begin{figure}
      \includegraphics[width=\textwidth]{going_back1.png}
    \end{figure}

\end{frame}


%%%%%%%%%%%%%%%%%%%%%%%%%%%%%%%%%%%%
\begin{frame}<beamer>{Going back to the original image size}

    \begin{figure}
      \includegraphics[width=\textwidth]{going_back2.png}
    \end{figure}

\end{frame}


%%%%%%%%%%%%%%%%%%%%%%%%%%%%%%%%%%%%
\begin{frame}{Going back to the original image size}

    \begin{figure}
      \includegraphics[width=\textwidth]{going_back.png}
    \end{figure}

\end{frame}


%%%%%%%%%%%%%%%%%%%%%%%%%%%%%%%%%%%%
\begin{frame}{Upsampling techniques}


\begin{itemize}
\item Replication
\item Transposed convolution
\item Pooling index memorization
\end{itemize}

\end{frame}

%%%%%%%%%%%%%%%%%%%%%%%%%%%%%%%%%%%%
\begin{frame}{Upsampling through replication}

  \begin{figure}
      \includegraphics[height=0.8\textheight]{upsampling.png}
    \end{figure}

\end{frame}


%%%%%%%%%%%%%%%%%%%%%%%%%%%%%%%%%%%%
\begin{frame}{Transposed convolution}

  \begin{columns}

    \begin{column}<1->{0.3\textwidth}
      \begin{center}
        \includegraphics[width=0.74\textwidth]{convolutional_layer}
        \\ \scriptsize{Convolutional layer}
      \end{center}
    \end{column}

    \begin{column}<2->{0.3\textwidth}
      \begin{center}
        \includegraphics[width=0.74\textwidth]{transposed_conv.png}
        \\ \scriptsize{Transposed convolutional layer}
      \end{center}
    \end{column}

    \begin{column}<3->{0.3\textwidth}
      \begin{center}
        \includegraphics[width=0.74\textwidth]{transposed_conv_stride.png}
        \\ \scriptsize{Transposed convolutional layer with stride $=2$}
      \end{center}
    \end{column}

  \end{columns}


\end{frame}

%% Je pense que cet example est avec stride =1 et on utilise stride = 2 pour l'upsampling
%% \begin{frame}{Transposed convolution: example}

%% \begin{figure}[ht]
%%   \centering
%%   \includegraphics[width=\textwidth]{TransposedConv_Example.pdf}
%% \end{figure}

%% \end{frame}


%%%%%%%%%%%%%%%%%%%%%%%%%%%%%%%%%%%%
%% \begin{frame}{Pooling index memorization}

%% \begin{figure}[ht]
%%   \centering
%%   \includegraphics[width=\textwidth]{index_unpooling}
%% \source{Divyanshu Mishra, Transposed convolution demystified, \url{http://towardsdatascience.com}}
%% \end{figure}


%% \end{frame}



%%%%%%%%%%%%%%%%%%%%%%%%%%%%%%%%%%%%%%%%%%%%%%%%%%%%%
\subsection{Receptive field}
%%%%%%%%%%%%%%%%%%%%%%%%%%%%%%%%%%%%%%%%%%%%%%%%%%
\frame{
  \frametitle{Receptive field}

  \begin{block}{Definition: links between neurons}
    In a NN, we say that neuron $a$ is linked to neuron $b$ if there is an oriented path in the corresponding graph going from $a$ to $b$.
  \end{block}


  \begin{block}{Definition}
    The \alert{receptive field} of a neuron in a NN is the set of \emph{input neurons} that are linked to that neuron.

    The size of the receptive field is an essential property when designing a fully-convolutional NN architecture.
  \end{block}

Note that in most cases, receptive fields are square shaped (border effects ignored).

}

%%%%%%%%%%%%%%%%%%%%%%%%%%%%%%%%%%%%
\begin{frame}<beamer>{Illustration}

  \begin{figure}
    \includegraphics[width=0.6\textwidth]{receptive_field.png}
  \end{figure}

\end{frame}

%%%%%%%%%%%%%%%%%%%%%%%%%%%%%%%%%%%%
\begin{frame}<beamer>{Illustration}

  \begin{figure}
  \includegraphics[width=0.6\textwidth]{receptive_field2.png}
  \end{figure}

\end{frame}

%%%%%%%%%%%%%%%%%%%%%%%%%%%%%%%%%%%%
\begin{frame}{Illustration}

  \begin{figure}
  \includegraphics[width=0.6\textwidth]{receptive_field3.png}
  \end{figure}

\end{frame}


%%%%%%%%%%%%%%%%%%%%%%%%%%%%%%%%%%%%
%% \begin{frame}{Receptive field evolution through an upsampling layer}

%%   \begin{figure}
%%   \includegraphics[width=0.8\textwidth]{receptive_field_upsampling.png}
%%   \end{figure}

%% \end{frame}

%%%%%%%%%%%%%%%%%%%%%%%%%%%%%%%%%%%%%%%%%%%%%%%%%%
\frame{
  \frametitle{Example}

  \begin{block}{}
    \begin{figure}
      \includegraphics[width=0.6\textwidth]{plain_convnet.png}
    \end{figure}
  \end{block}


  \begin{block}{What is the width of the receptive field of the neurons in the last layer?}%quizzblock
    \begin{enumerate}
    \item $n$
    \item $1 + 2 \times n$
    \item $1 + 2 \times (n+1)$
    \item It depends on the input image
    \end{enumerate}

  \end{block}%quizzblock

  \mode<beamer>{
  \pause

  Answer: $1 + 2 \times (n+1)$
  }
  }
%%%%%%%%%%%%%%%%%%%%%%%%%%%%%%%%%%%%%%%%%%%%%%%%%%
%% \frame<handout>{
%%   \frametitle{Example}

%%   \begin{block}{What is the width of the receptive field of the neurons in the last layer?}%quizzblock
%%     \begin{figure}
%%       \includegraphics[width=\textwidth]{plain_convnet.png}
%%     \end{figure}
%%   \end{block}%quizzblock
%% }

%%%%%%%%%%%%%%%%%%%%%%%%%%%%%%%%%%%%
\begin{frame}{Effective receptive field: intuition}

  \begin{figure}
    \includegraphics[width=0.6\textwidth]{receptive_field3.png}
  \end{figure}

\end{frame}


%%%%%%%%%%%%%%%%%%%%%%%%%%%%%%%%%%%%

\begin{frame}{Conclusion on the receptive field}

  When choosing or designing a NN architecture:
\begin{itemize}
\item Establish its minimal required size with respect to the application
\item Make it larger than that...
\end{itemize}

\end{frame}



\section{U-net architecture}

%%%%%%%%%%%%%%%%%%%%%%%%%%%%%%%%%%%%%%%%%%%%%%%%%%%%%%%%%%%%
\begin{frame}{U-net architecture}

  %..........
  \begin{block}{Architecture}
  \begin{itemize}{}
  \item Fully convolutional network, with an encoder and a decoder.

  \item The \textbf{encoder} (contracting path) is inspired by typical classification networks:
    %..............................
    \begin{itemize}{}
    \item    a series basic blocks: Convolution (3x3), Activation (RELU) and max-pooling 2x2 for downsampling.
  \item At each downsampling step, the number of channels (features) is multiplied by 2.
    \end{itemize}
    %..............................
  \item The \textbf{decoder} (expanding path)
    %..............................
    \begin{itemize}{}
    \item upsampling the feature map
    \item concatenate corresponding features from contracting path $\Rightarrow$   \textcolor{blue}{\underline{\textbf{ Skip connection}}}
    \item Basic block: Convolution (3x3), Activation (RELU)
    \end{itemize}
    %..............................
  \item A final Convolution 1x1 with the output number of classes.
  \end{itemize}
  %..............................

    \end{block}
  %..............................

{\small    Ronneberger, O., Fischer, P., and Brox, T. (\textbf{2015}). U-net: Convolutional networks for biomedical image segmentation. MICCAI 2015.} University of Freiburg, Germany.

  \end{frame}

\begin{frame}{Skip connection}
  %..........
  \begin{block}{The key idea}

  Skip connection combine:
  %..............................
  \begin{itemize}{}
  \item  spatial information from the encoding part
  \item  with up-sampled deeper layers with more relevant feature representations.
  \end{itemize}
  %..............................

    \end{block}

\end{frame}

%%%%%%%%%%%%%%%%%%%%%%%%%%%%%%%%%%%%%%%%%%%%%%%%%%%%%%%%%%%%
\begin{frame}{U-net architecture}
\centerline{  \includegraphics[width=0.8\textwidth]{../images/unet_architecture.png}}

  \end{frame}
%%%%%%%%%%%%%%%%%%%%%%%%%%%%%%%%%%%%%%%%%%%%%%%%%%%%%%%%%%%%

\begin{frame}[fragile]{U-net implementation (keras)}
  \begin{small}
    %\begin{Verbatim}[commandchars=\\\{\}]
    \begin{alltt}
  def \textcolor{blue}{modular_u_net}(shape, nb_levels = 4, nb_filters=32,output_channels=1):
    x = Input(shape); first_layer=[]; first_layer.append(x); skip_layers=[]
    \textcolor{blue}{# ENCODER}
    for lev in range(nb_levels):
        x = unet_block(nb_filters, res, "conv"+str(lev+1), batch_norm)(x)
        if(lev < nb_levels -1):#not the last one
            skip_layers.append(x)
            x = unet_down(conv_sampling, nb_filters)(x)
            nb_filters = nb_filters * 2
    \textcolor{blue}{# DECODER}
    for lev in range(nb_levels-1):
        x = unet_up(conv_sampling, nb_filters)(x)
        \textcolor{blue}{\textbf{x = concatenate([x, skip_layers.pop()], axis=3)}}
        nb_filters= int(nb_filters/2)
        x = unet_block(nb_filters, res, "conv"+str(nb_levels+lev+1), batch_norm)(x)
    xFinal = Conv2D(output_channels, 1, activation="sigmoid", name="out")(x)
    return Model(first_layer.pop(), xFinal)
\end{alltt}
  \end{small}
  \end{frame}

%%%%%%%%%%%%%%%%%%%%%%%%%%%%%%%%%%%%%%%%%%%%%%%%%%%%%%%%%%%%
\begin{frame}[fragile]{U-net block}

  %..........
  \begin{alltt}
    def \textcolor{blue}{unet_block}(nb_fil, res, name, batch_norm):
    def fn(tensor):
        x = Conv2D(nb_fil, 3, padding="same", name="\%s_1" \% name)(tensor)
        if batch_norm:
            x = BatchNormalization()(x)
        x = Activation("relu")(x)
        x = Conv2D(nb_fil, 3, padding="same", name="\%s_2" \% name)(x)
        if batch_norm:
            x = BatchNormalization()(x)
        if res:
            skip=Conv2D(nb_fil,3,padding="same",name="\%s_res"\% name)(tensor)
            x = add([x, skip])
        x = Activation("relu")(x)
        return x
    return fn
    \end{alltt}
  \end{frame}
%%%%%%%%%%%%%%%%%%%%%%%%%%%%%%%%%%%%%%%%%%%%%%%%%%%%%%%%%%%%
\begin{frame}[fragile]{U-net down block}
  %..........
  \begin{alltt}
def \textcolor{blue}{unet_down}(conv_sampling, nb_filters=None):
    if conv_sampling:
        if nb_filters is None:
            raise ValueError("When conv_sampling is True, nb_filters should be given")
        def fn(tensor):
            x = Conv2D(nb_filters, 3, padding="same", strides=(2, 2))(tensor)
            return x
    else:
        def fn(tensor):
            x = MaxPooling2D(pool_size=(2, 2))(tensor)
            return x
    return fn

    \end{alltt}
  \end{frame}
%%%%%%%%%%%%%%%%%%%%%%%%%%%%%%%%%%%%%%%%%%%%%%%%%%%%%%%%%%%%
\begin{frame}[fragile]{U-net up block}
  %..........
  \begin{alltt}
def \textcolor{blue}{unet_up}(conv_sampling, nb_filters):
    if conv_sampling:
        if nb_filters is None:
            raise ValueError("When conv_sampling is True, nb_filters should be given")
        def fn(tensor):
            x = Conv2DTranspose(nb_filters, 3, padding="same", strides=(2, 2))(tensor)
            return x
    else:
        def fn(tensor):
            x = UpSampling2D(size=(2, 2))(tensor)
            return x
    return fn

    \end{alltt}
  \end{frame}

\section{Unet variants}

%%%%%%%%%%%%%%%%%%%%%%%%%%%%%%%%%%%%%%%%%%%%%%%%%%%%%%%%%%%%
\begin{frame}{Unet variants}
  %..............................
  \begin{itemize}{}
  \item Add residual connections in the convolutional blocks
  \item Unet ++: A nested u-net architecture for medical image segmentation. 2018.
  \item Attention u-net
  \item UnetR: U-net Transformer
  \item Sharp U-net
  \item Motion U-net
  \item ...

  \end{itemize}
  %..............................

  Siddique, N., Paheding, S., Elkin, C. P., \& Devabhaktuni, V. (2021). U-net and its variants for medical image segmentation: A review of theory and applications. IEEE Access, 9.

\end{frame}
%%%%%%%%%%%%%%%%%%%%%%%%%%%%%%%%%%%%%%%%%%%%%%%%%%%%%%%%%%%%
\begin{frame}{Unet++}
  %..........
  \begin{block}{Dense skip pathways}
    %..............................
    \begin{itemize}{}
    \item  Aim: reduce semantic gap between the encoder and decoder feature maps.
    \item $x^{0,1} = H[x^{0,0},U(x^{1,0})]$; $x^{0,2} = H[x^{0,0},x^{0,1},U(x^{1,1})]$; $x^{0,2} = H[x^{0,0},x^{0,1},x^{0,2},U(x^{1,2})]$...

    \end{itemize}
    %..............................

    \end{block}

  %::::::::::::::::::::::::::::::
  \begin{columns}
    \begin{column}{0.65\textwidth}
  \centerline{        \includegraphics[width=0.9\textwidth]{../images/UnetPlusPlus_01.png}}
    \end{column}
    \begin{column}{0.35\textwidth}
        %Rahman Siddiquee, M. M., Tajbakhsh, N., \& Liang, J
  \begin{footnotesize}  Zhou, Z. et al. (2018). Unet++: A nested u-net architecture for medical image segmentation. In DL in medical image analysis and multimodal learning for clinical decision support (pp. 3-11).
    \end{footnotesize}

    \end{column}
  \end{columns}
  %::::::::::::::::::::::::::::::


  \end{frame}
%%%%%%%%%%%%%%%%%%%%%%%%%%%%%%%%%%%%%%%%%%%%%%%%%%%%%%%%%%%%
\begin{frame}{Unet++}
  %::::::::::::::::::::::::::::::
  \begin{columns}
    \begin{column}{0.65\textwidth}
        \includegraphics[width=0.8\textwidth]{../images/UnetPlusPlus_01.png}

        \includegraphics[width=0.8\textwidth]{../images/UnetPlusPlus_02.png}

    \end{column}
    \begin{column}{0.35\textwidth}
      \includegraphics[height=0.95\textheight]{../images/UnetPlusPlus_03.png}
    \end{column}
  \end{columns}
  %::::::::::::::::::::::::::::::
\end{frame}

%%%%%%%%%%%%%%%%%%%%%%%%%%%%%%%%%%%%%%%%%%%%%%%%%%%%%%%%%%%%
\begin{frame}{Attention U-net}
\centerline{  \includegraphics[width=0.8\textwidth]{../images/unet_attention_architecture.png}}

  \end{frame}
%%%%%%%%%%%%%%%%%%%%%%%%%%%%%%%%%%%%%%%%%%%%%%%%%%%%%%%%%%%%
\begin{frame}{Attention u-net}
%https://www.youtube.com/watch?v=KOF38xAvo8I

  %..............................
  \begin{itemize}{}
  \item  Highlights relevant activations during training
  \item   Reduces computational ressources on irrelevant activations
  \item   Better generalization
  \end{itemize}
  %..............................

  %..........
  %% \begin{block}{Hard attention}
  %%   %..............................
  %%   \begin{itemize}{}
  %%   \item    One region of an image at a time
  %%   \item Relevant regions are cropped
  %%   \item     Network can pay attention or not, binary choice
  %%   \item not differentiable, Backpropagation cannot be used, needs reinforcement learning)
  %%   \end{itemize}
  %%   %..............................
  %%   \end{block}


    %..........
    \begin{block}{Soft attention}
      %..............................
      \begin{itemize}{}
      \item Attention weights the relevance of different parts of the image
      \item The more relevant, the larger weight
      \item Backpropagation can be used
      \item The weights are trained
      \item Adds weights to pixels based on their relevance

      \end{itemize}
      %..............................

      \end{block}


  Oktay, O., Schlemper, J., Folgoc, L. L., Lee, M., Heinrich, M., Misawa, K., ... \& Rueckert, D. (2018). Attention u-net: Learning where to look for the pancreas. arXiv preprint arXiv:1804.03999.

\end{frame}


%%%%%%%%%%%%%%%%%%%%%%%%%%%%%%%%%%%%%%%%%%%%%%%%%%%%%%%%%%%%
\begin{frame}{Attention gate}

  %..............................
  \begin{itemize}{}
  \item Takes two inputs:
    %..............................
    \begin{itemize}{}
    \item x from the skip connection, with better spatial information
    \item g from the lower layer, with better feature representation
    \end{itemize}
    %..............................



  \end{itemize}
  %..............................
\centerline{  \includegraphics[width=0.8\textwidth]{../images/unet_attention_gate.png}}
\end{frame}


%%%%%%%%%%%%%%%%%%%%%%%%%%%%%%%%%%%%%%%%%%%%%%%%%%%%%%%%%%%%
\begin{frame}[fragile]{Attention gate code}
%  \centerline{  \includegraphics[width=0.5\textwidth]{../images/unet_attention_gate.png}}
  \begin{small}
%https://www.youtube.com/watch?v=KOF38xAvo8I
\begin{verbatim}
def attention_block_2d(x, g, inter_channel):
    # x is the skip connexion
    x2 = Conv2D(inter_channel, [1, 1], strides=[1, 1])(x)
    # g is the up-sampled deeper level feature (stride is then 1)
    g2 = Conv2D(inter_channel, [1, 1], strides=[1, 1])(g)

    # add both branches and then activation
    f = Activation('relu')(add([x2, g2]))

    # get the scalar value, then apply sigmoid
    psi_f = Conv2D(1, [1, 1], strides=[1, 1])(f)
    rate = Activation('sigmoid')(psi_f)

    # multiply
    att_x = multiply([x, rate])
    return att_x


\end{verbatim}
\end{small}


%https://github.com/lixiaolei1982/Keras-Implementation-of-U-Net-R2U-Net-Attention-U-Net-Attention-R2U-Net.-/blob/master/network.py


\end{frame}


\section{Google Colab}

%%%%%%%%%%%%%%%%%%%%%%%%%%%%%%%%%%%%%%%%%%%%%%%%%%%%%%%%%%%%
\begin{frame}{Google Colab}
  %..............................
  \begin{itemize}{}
  \item Go to the website https://colab.research.google.com/
  \item Open a new notebook or an existing one
  \item For existing notebooks, upload it to your google drive
  \item Copy it in the colab content. For this purpose:
    %..............................

\textcolor{blue}{     !cp /content/gdrive/MyDrive/Colab\ Notebooks/my\_notebook.ipynb /content"}

    %..............................
  \item Copy additional files:
    %..............................

\textcolor{blue}{\small{     !cp /content/gdrive/MyDrive/Colab\ Notebooks/Unet\_MAREVA/parameters.py /content}}

\textcolor{blue}{\small{      !cp -r /content/gdrive/MyDrive/Colab\ Notebooks/Unet\_MAREVA/dl\_tools /content}}

%..............................
{\bf{Replace "Unet\_MAREVA" with your directory name on google drive.}}


  \item In the first cell of your notebook add the following two lines:

\textcolor{blue}{    from google.colab import drive}

\textcolor{blue}{    drive.mount('/content/gdrive')}
  \end{itemize}
  %..............................

\end{frame}

\section{Conclusions}
%%%%%%%%%%%%%%%%%%%%%%%%%%%%%%%%%%%%%%%%%%%%%%%%%%%%%%%%%%%%
\begin{frame}{Conclusions}
%..............................
  \begin{itemize}{}
\item U-net and its variants continue to be a prominent method for semantic segmentation
\item For medical images and other domains
\item Rmsprop or Adam can be used as default optimizers
\item Jaccard2 or Cross-Entropy can be used as a baseline
\item Loss is computed pixel-wise
\item Working on specific losses enforcing structure is an interesting research topic

\end{itemize}
%..............................

  \end{frame}


\begin{frame}{Residual learning (ResNet)}

  Shortcut connection: Allow deeper networks with less risk of performance loss.
  %::::::::::::::::::::::::::::::
  \begin{columns}
    \begin{column}{0.2\textwidth}


 \centerline{     Residual block}
    \end{column}
    \begin{column}{0.8\textwidth}
      \centerline{  \includegraphics[width=0.95\textwidth]{../images/resnet_block.png}}
    \end{column}
  \end{columns}
  %::::::::::::::::::::::::::::::


\small{  He, Kaiming, et al. "Deep residual learning for image recognition." Proceedings of the IEEE conference on computer vision and pattern recognition. 2016.}

\small{\href{https://openaccess.thecvf.com/content_cvpr_2016/papers/He_Deep_Residual_Learning_CVPR_2016_paper.pdf}{https://openaccess.thecvf.com/content\_cvpr\_2016/papers/He\_Deep\_Residual\_Learning\_CVPR\_2016\_paper.pdf}}
\end{frame}
%% %%%%%%%%%%%%%%%%%%%%%%%%%%%%%%%%%%%%%%%%%%%%%%%%%%%%%%%%%%%%
%% \begin{frame}{Res-Unet}
%%     %::::::::::::::::::::::::::::::
%%   \begin{columns}
%%     \begin{column}{0.5\textwidth}


%%  \centerline{\small{\href{https://ieeexplore.ieee.org/document/8959021}{  \parbox[c]{0.9\textwidth}{
%% Zhang, Z., Liu, Q., \& Wang, Y. (2018). Road extraction by deep residual u-net. IEEE Geoscience and Remote Sensing Letters, 15(5), 749-753.}}}
%% }
%%     \end{column}
%%     \begin{column}{0.4\textwidth}
%%       \centerline{  \includegraphics[width=0.80\textwidth]{../images/resunet.png}}
%%     \end{column}
%%   \end{columns}
%%   %::::::::::::::::::::::::::::::

%%   \end{frame}
%% %%%%%%%%%%%%%%%%%%%%%%%%%%%%%%%%%%%%%%%%%%%%%%%%%%%%%%%%%%%%
%% \begin{frame}{Res-Unet ++}
%%     %::::::::::::::::::::::::::::::
%%   \begin{columns}
%%     \begin{column}{0.5\textwidth}


%%  \centerline{\small{\href{https://ieeexplore.ieee.org/document/8959021}{\parbox[c]{0.9\textwidth}{Jha, D., Smedsrud, P. H., Riegler, M. A., Johansen, D., De Lange, T., Halvorsen, P., \& Johansen, H. D. (2019, December). Resunet++: An advanced architecture for medical image segmentation. In 2019 IEEE international symposium on multimedia (ISM) (pp. 225-2255). IEEE.}}}
%% }
%%     \end{column}
%%     \begin{column}{0.4\textwidth}
%%       \centerline{  \includegraphics[width=0.95\textwidth]{../images/resunet++.png}}
%%     \end{column}
%%   \end{columns}
%%   %::::::::::::::::::::::::::::::

%%   \end{frame}
%%%%%%%%%%%%%%%%%%%%%%%%%%%%%%%%%%%%%%%%%%%%%%%%%%%%%%%%%%%%
\begin{frame}{Google colab}
  %..............................
  \begin{itemize}{}
  \item Upload files (notebook, dl\_tools directory and parameters.py file) to your google drive
  \item mount your google drive
    \begin{itemize}
\item
  from google.colab import drive
\item
  drive.mount('/content/gdrive')
      \end{itemize}

  \item change directory
    \begin{itemize}
\item
import os
\item
  os.chdir("gdrive/My Drive/")
      \end{itemize}

  \end{itemize}
  %..............................

  \end{frame}
%%%%%%%%%%%%%%%%%%%%%%%%%%%%%%%%%%%%%%%%%%%%%%%%%%%%%%%%%%%%
%% \begin{frame}[fragile]{Dashboard}
%% \begin{alltt}
%% tensorboard --logdir=runs   (path of logfiles) runs directory

%% http://localhost:6006/

%% from torch.utils.tensorboard import SummaryWriter

%% writer = SummaryWriter("runs/mnist")

%% for n_iter in range(100):
%%     writer.add_scalar('Loss/train', np.random.random(), n_iter)
%%     writer.add_scalar('Loss/test', np.random.random(), n_iter)
%%     writer.add_scalar('Accuracy/train', np.random.random(), n_iter)
%%     writer.add_scalar('Accuracy/test', np.random.random(), n_iter)

%% img_grid = torchvision.utils.make_grid(example_data)
%% writer.add_image('mnist_images',image_grid)
%% writer.close()


%% logdir = "logs/scalars/" + datetime.now().strftime("%Y%m%d-%H%M%S")
%% file_writer = tf.summary.create_file_writer(logdir + "/metrics")
%% file_writer.set_as_default()

%% tf.summary.scalar('learning rate', data=learning_rate, step=epoch)

%% \end{alltt}
%% \end{frame}
%%%%%%%%%%%%%%%%%%%%%%%%%%%%%%%%%%%%
%%%%%%%%%%%%%%%%%%%%%%%%%%%%%%%%%%%%
%\input{efficientNet}
%%%%%%%%%%%%%%%%%%%%%%%%%%%%%%%%%%%%
%%%%%%%%%%%%%%%%%%%%%%%%%%%%%%%%%%%%


%%%%%%%%%%%%%%%%%%%%%%%%%%%%%%%%%%%%%%%%%%%%%%%%%%
%%%%%%%%%%%%%%%%%%%%%%%%%%%%%%%%%%%%%%%%%%%%%%%%%%
%%%%%%%%%%%%%%%%%%%%%%%%%%%%%%%%%%%%%%%%%%%%%%%%%%
% FROM ETIENNE's slides
%\input{attention}
%\input{attention_towardsai}
%%%%%%%%%%%%%%%%%%%%%%%%%%%%%%%%%%%%%%%%%%%%%%%%%%
%%%%%%%%%%%%%%%%%%%%%%%%%%%%%%%%%%%%%%%%%%%%%%%%%%
%%%%%%%%%%%%%%%%%%%%%%%%%%%%%%%%%%%%%%%%%%%%%%%%%%
\section*{References}

%\bibliographystyle{elsart-num}
%\bibliographystyle{plain}
%\bibliography{D:/marcoteg-svn/papers/HDR/HDR_biblio}
\frame[allowframebreaks]{

  \scriptsize

  \frametitle{References}

  %\bibliographystyle{amsalpha}
  %\bibliographystyle{apalike}
\bibliographystyle{plain}
  \bibliography{edf.bib}

  \normalsize

}

\end{document}


%%%%%%%%%%%%%%%%%%%%%%%%%%%%%%%%%%%%%%%%%%%%%%%%%%%%%%%%%%%%
